\documentclass{article}
\PassOptionsToPackage{compress}{natbib}
\usepackage[dblblindworkshop]{neurips_2025}
% \usepackage[dblblindworkshop, final]{neurips_2025}
\usepackage{graphicx} % Required for inserting images

\title{SMM}
\author{Roussel Rahman}
\author{Jeff Shrager}
\date{September 2025}

% Optional math and formatting packages
\usepackage{amsmath, amssymb}
\usepackage{graphicx}
\usepackage{booktabs}
\usepackage{xcolor}

\newcommand{\rrsuggestion}[1]{\textcolor{blue}{#1}}

% RR: Setting Boxes vs Colored texts for links, citations etc.

\usepackage{hyperref}
%\AtBeginDocument{
\hypersetup{
    colorlinks=true,
    linkcolor=blue,
    citecolor=teal, % ForestGreen, MidnightBlue, Magenta
    filecolor=green,
    urlcolor=blue}

\title{A Small Math Model: Reimplementing Siegler and Shrager's Addition Model in an LLM-Inspired Architecture}

% Alternate titles:

% \title{A Small Math Model: Recreating a Cognitive Model of Learning Addition Problems in an LLM-Inspired Architecture}

% \title{A Small Math Model: Teaching a Transformer-based Neural Network how to Add like Children}

% \title{A Small Math Model: Enulating Children/Human Learning in Addition Problems in an LLM-Inspired Architecture}

\author{
  Jeff Shrager\\
  Stanford University\\
  Stanford, CA 94305\\
  \texttt{jshrager@gmail.com}
  \And
  Roussel Rahman\\
  Stanford University\\
  Stanford, CA 94305\\
  \texttt{roussel.rahman@gmail.com}
}

\begin{document}

\maketitle

\begin{abstract}
\citeauthor{siegler1984strategychoices}'s influential 1984 strategy choice model (S\&S84) anticipated principles now central to Large Language Models (LLMs), including confidence-based retrieval, backup strategies, and learning from naturalistic data. We present the ``Small Math Model'' (SMM), a modern neural implementation of S\&S84 extended with a preceding learning-to-count stage. Like LLMs, the SMM processes symbolic sequences and produces probabilistic outputs, reproducing hallmark behaviors predicted by strategy choice theory. While pedagogically useful, the SMM’s primary contributions are: (a) modernizing the Siegler program within an LLM-inspired framework, showing that a single mechanism can account for diverse developmental patterns, and (b) providing a platform for investigating number common sense in AI. Future work introduces multiple strategies and learned embeddings to test whether the model can capture emergent number properties and learn adaptive strategy use.
\end{abstract}


\section{Introduction}
The development of mathematical cognition has long been understood through the lens of strategy choice theory \citep{siegler1984strategychoices, siegler1987perils, siegler1991microgenetic, shrager1998scads}. Children use multiple strategies---counting, retrieval, decomposition---and select among them probabilistically based on confidence and associative strength. \citeauthor{siegler1984strategychoices}'s framework emphasized mechanisms such as backup strategies and a ``feeling of knowing,'' concepts that align with probabilistic inference, uncertainty estimation, and fallback mechanisms in modern AI.

Meanwhile, Large Language Models (LLMs) have revolutionized natural language processing by scaling these same principles. LLMs produce probabilistic outputs, estimate confidence via entropy, and dynamically select strategies from massive learned corpora. This parallel motivates a re-examination of early cognitive frameworks through the lens of modern architectures.

We present the Small Math Model (SMM), a shallow (three-layer) neural network designed to explore the counting–addition conflict central to early math development. The SMM allows us to revisit classical questions in a tractable contemporary architecture, and leverage contemporary analytic tools. Although the SMM may have pedagogical applications, our primary goals are twofold: (a) to bring the Siegler modeling program into the LLM era, demonstrating that a single mechanism can still account for diverse developmental patterns within a modern framework, and (b) to provide a platform for studying number common sense in LLMs, which has been shown to be limited in other contexts. 

This way, the SMM serves both as a modern continuation of Siegler and Shrager’s tradition and as a testbed for probing fundamental issues in AI reasoning. Our immediate next steps involve expanding beyond a single strategy, allowing the model to learn when to apply different strategies, and incorporating number embeddings to test whether they can capture emergent numerical properties and support improved strategy selection. 

\section{SMM $<->$ LMM Analysis}

\section{Architecture}

The Small Math Model (SMM) represents numbers using learned embeddings. Each integer is mapped to a dense vector in a continuous space, enabling the model to capture emergent relationships among numbers (for example, even versus odd, relative magnitude, or common configurations such as base-five hand groupings). These embeddings feed into a recurrent processing core that implements a mixture of symbolic strategies, including finger counting and direct addition.

A central gating mechanism allocates probability mass across strategies based on past performance and contextual confidence, providing a computational analogue of adaptive strategy selection in human learners. In addition, the model incorporates an \emph{attention gate} that learns to modulate which
input arguments are emphasized in different contexts. Through training, the model acquires patterns of attention that support distinct strategies: in counting, attention tends to amplify the representation of the current number in a sequence, facilitating stepwise incrementing; in addition, attention tends to highlight both operands simultaneously, capturing their relational structure. These attentional allocations are not hand-coded but emerge as a consequence of learning. To avoid unstable behavior early in training, the attention gate is held off for an initial warm-up period; once enabled, it allows the model to selectively amplify relevant signals and suppress distractors, sharpening the competition among strategies.

A distinctive feature of the SMM is its incorporation of explicit \emph{finger counting} as an auxiliary reasoning strategy. Confidence is estimated from the entropy of the output distribution; when confidence falls below a threshold, the model defers to an external finger-counting procedure that generates the correct result step by step. The output of this routine is then fed back into training, analogous to the way LLMs can be augmented with external tools or chains-of-thought and subsequently distilled back into the base model.

Counting and addition are implemented as separable but overlapping modules. Counting is trained from the very beginning of each run, while addition can be introduced at different timepoints to investigate how prior knowledge influences subsequent learning. After each problem, strategy-specific confidence values are updated according to correctness, and these values govern stochastic selection of strategies on subsequent trials. This architecture
therefore supports mechanistic analysis of how new procedures compete with, reinforce, or interfere with existing ones in the course of skill acquisition.

Training follows a curriculum schedule in which problem difficulty and task type (counting vs. addition) are varied according to Gaussian flows. This is directly inspired by curriculum learning practices in LLM training, where data distributions are staged over time to improve stability and generalization. Learning rates and confidence thresholds are annealed across steps, further paralleling hyperparameter scheduling in large-scale models.

In sum, the SMM modernizes and extends earlier work by Shrager et al.\ by explicitly drawing analogies to contemporary LLM architectures: embeddings, gating/attention, cross-entropy loss, curriculum learning, and tool use. This vastly simplified architecture provides a testbed for studying how counting priors, confidence estimation, and quasi-attention interact in a controlled setting, with the ultimate goal of bridging symbolic and neural perspectives on mathematical reasoning.


\section{Methods}
The SMM processes input sentences of the form ``a operator b'' where $a,b \in \{0,\ldots,9\}$ and operator $\in \{+\}$. Inputs are encoded as one-hot vectors and passed through a three-layer feedforward network trained with cross-entropy loss to predict the sum distribution.

\paragraph{Curriculum Learning.} Training follows a Gaussian curriculum schedule that gradually increases task difficulty. Early in training, simpler problems (e.g., sums with smaller operands) dominate; over time, harder problems appear with greater probability. This smooth progression mirrors natural learning trajectories while avoiding abrupt jumps.

\paragraph{Logging and Evaluation.} The implementation includes structured logging of accuracy, entropy, and error distributions across problem types. These diagnostics allow comparison to known developmental error patterns, such as over-reliance on counting or failure to generalize across contexts.

\section{Results}

To probe the dynamics of interference, we conducted a sweep varying the point at which
addition training was introduced (2000, 4000, …, 14000 steps). Counting was trained from the
very beginning in all cases. Figure~\ref{fig:addition_interference} shows the moving-window
correctness (2000-sample average) on addition problems as a function of training step.

The results reveal both constructive and destructive interference between counting and
addition. When addition is introduced very early (e.g., 2000 steps), learning is unstable and
plateaus at moderate accuracy (around 0.5). When introduced at midpoints (6000–8000 steps),
addition learning benefits from prior counting knowledge, rising sharply to high accuracy
levels (0.7+). In contrast, when addition is introduced very late (12000–14000 steps),
correctness remains suppressed, rarely exceeding 0.2. Thus, prior experience with counting
can either scaffold or hinder the acquisition of addition, depending critically on the timing
of introduction.

These findings extend Siegler and Shrager’s framework into a computational setting,
demonstrating how adaptive strategy competition can give rise to nontrivial transfer effects.
The embedding-based architecture further allows us to investigate whether number embeddings
capture meaningful latent properties (e.g., even/odd or base-five hand configurations), which
we leave to future work.

\begin{figure}[t]
  \centering
  \includegraphics[width=0.8\linewidth]{addition_start_step_add_correctness2000.png}
  \caption{Addition correctness as a function of training step, for different
  addition-start times. Counting is always trained from the beginning. Constructive
  and destructive interference emerge depending on when addition is introduced.}
  \label{fig:addition_interference}
\end{figure}
\section{Discussion and Next Steps}
The SMM illustrates how modern neural implementations can instantiate and extend classical cognitive theories. By reframing Siegler and Shrager's ideas in computational terms aligned with LLMs, we obtain a platform for mechanistic analysis of symbolic reasoning in tractable domains, opportunities to probe representations directly (e.g., via similarity analysis or attention visualization), and a bridge between cognitive science and AI, highlighting shared principles of probabilistic selection, confidence estimation, and adaptive strategy use.

Although the SMM may have pedagogical utility, its main contributions are twofold: (a) it brings the Siegler modeling program into the LLM era, demonstrating that a single mechanism can reproduce diverse developmental patterns within a contemporary architecture; and (b) it provides a platform for investigating number common sense in systems modeled after LLMs, an ability that remains fragile in current large-scale AI.

Looking forward, the most important step is to move beyond a single strategy. We plan to implement multiple strategies and allow the model to learn which are appropriate for which types of problems. This expansion will allow closer alignment with empirical findings on children’s adaptive strategy choice, while testing whether embeddings can capture properties of numbers (e.g., parity, magnitude relations, proximity to 5 or 10) that influence strategy efficiency. 

While the current implementation uses one-hot encodings, we will introduce \textbf{learned embeddings} for numbers and operators. Embeddings should enable the discovery of emergent number properties and support strategy-specific generalizations. To test this, we will design pseudo-strategies with error profiles dependent on number features---for instance, a ``finger counting'' strategy that excels at $+1$ problems, or a $+5$ strategy that works best when an operand is 5. By varying error proneness we can examine whether embeddings capture such latent features. Probes (e.g., linear classifiers on embeddings) and behavioral splits (e.g., 1+ problems vs. 5+ problems) will serve as evaluation tools.

Together, these extensions will provide a richer testbed for exploring how modern neural systems can internalize structured numerical knowledge, reconcile competing strategies, and develop more robust common sense. In doing so, the SMM continues Siegler’s tradition of mechanistic modeling while directly engaging with central challenges in today’s AI research.


\section*{\rrsuggestion{Inventory -- TBD}}



\subsection*{Some high-level points}

(RR: Some points I thought important in general. Some may be too much to include within four pages.)

Investigating reasoning through aggregate scores obscure the use of different strategies, but true/accurate explanations need to specify the exact strategies \citep{newell1958elements, simon1971human, simon1976substantive, siegler1987perils, shrager1998scads, gigerenzer2020explain}. Reasoning and problem-solving in complex environments require complex skill organization, and it may be too complex to identify the elements of reasoning. A promising approach: explain reasoning in complex environments through the changes in strategies (i.e., their use) as new strategies are learned, in the same vein as \citet{siegler1984strategychoices} and \citet{shrager1998scads}.

The SMM forms the foundation of such an experimental testbed. A NN trained to learn addition by first learning to count and then building upon this ability to learn to add.

Catastrophic Forgetting vs Grokking. (What is grokking? \citet{nanda2023progress}) Catastrophic forgetting may stem from strict adherence to one strategy in uncertain situations that demand flexible use of different strategies. For example, in solving addition problems from memory, humans may revert to counting when confidence is low in the answer recalled from memory.

(what we may show here) The changes in strategies captured through the changes in step distributions. For humans, changes in strategy usage are reflected in changes in RT distributions. Here, we emulate the changes in the SMM. (specifics in a sentence) We show that .... shows the promise of the "microgenesis" of mathematical reasoning in environments of controllable complexity as a path to explaining AI reasoning in complex environments, such as mathematical problem-solving.

\subsection*{Some relevant research}
Catastrophic forgetting a well-established phenomenon in NNs, first observed by \citet{mccloskey1989catastrophic}
``New learning may interfere catastrophically with old learning when networks are trained sequentially'' \citep[][pp. 109]{mccloskey1989catastrophic}. While (some explanations), a detailed understanding of the phenomenon remains elusive and the problem is further exacerbated as LLMs become more complex with finetuning for specific domains and enhancement with different mechanisms (e.g., code execution, web search, and reinforcement learning). Consequently, even though the LLMs frequently demonstrate catastrophic forgetting in complex environments \citep[][]{shojaee2025illusion, rahman2025llm}, their causes are hard to specify and address. The SMM emulates the catastrophic forgetting and provides a testbed for existing and new explainable AI (XAI) tools.

(The need to study strategies)
Grokking: A sudden, sharp rise in performance after prolonged period of training. Nanda et al explained grokking in modular arithmetic. explanation: the model shifts from memorization to generalization based on a learned pattern.

Reverse grokking, a lesser-studied but intriguing phenomenon, flips the script: a model generalizes well early in training (high validation accuracy) but, with further training, loses generalization, reverting to overfitting or poor validation performance while training accuracy remains high. It’s as if the model forgets its insights and gets stuck in a memorization rut.

Shojaee et al's demos show that learning general strategies is not guaranteed. The model performance catastrophically collapsed in the ToH task, whereas there exists multiple strategies that yield optimal performance (i.e., shortest number of moves to goal state).

Generally, In complex environments, human learning curves go through pdls due to strategic exploration and discovery. 


\subsection{About Results}
SMM: (1) Performance drop when learning to add after learning to count: emulate Reverse Grokking. (2) The subsequent performance leap Counting to Addition -- emulate grokking. Explanation: The model indeed forgets counting to some extent and their addition skills are also not strong enough early in learning, resulting in the drop. Afterwards, the model become proficient in learning and also learns its counting skills back from repeated trials, resulting in improvement.


\subsection{For future}

0. Changes in the Distributions of Associations

0a. plot the softmax distribution over answer space for the same problem in epochs far apart (1000?)

0b. A measure for the changes in distributions across trials. We can "center" each distribution around correct answer and then use KL divergence or other distributional measures. This one can be compared with accuracy curve. If they reflect each other (in other words, are highly correlated), that would indicate the distributional changes may provide able indicators of changes at one level lower of AI reasoning with learning. The SMM allows us to trace our way into lower levels, such as in embedding spaces and down to input level.


1. Mechanistic explanations of SMM -- that is, identifying the circuits corresponding to counting vs addition strategies.

2. Distances in Embedding Space for learned strategies

3. Activations based on inputs

Generally, this way, the SMM provides a foundation to develop, test, and validate new tools to explain AI reasoning in complex domains such as mathematical problem solving and beyond.

\bibliographystyle{apalike}
\bibliography{math_reasoning_bibliography}

\end{document}
